\chapter*{GENERAL INTRODUCTION}
\addcontentsline{toc}{chapter}{GENERAL INTRODUCTION}
\markboth{GENERAL INTRODUCTION}{}

Modern organizations continuously generate customer-related data across heterogeneous sources: transactional systems, e-commerce platforms, marketing tools, behavioral event tracking, and operational applications. While the volume and variety of this information can enable more personalized experiences and better decision-making, it also creates a major engineering challenge: raw data arrives with different formats, update frequencies, quality levels, and semantics, which makes end-to-end analytics difficult without a structured data platform.

In many environments, customer data pipelines are implemented as a collection of scripts and manual exports. This approach does not scale when the number of sources increases, when data volume grows, or when analytics teams require frequent refreshes. The absence of standardized ingestion, versioned storage, and automated validation often leads to recurring issues such as duplicated work, inconsistent dashboards, and low trust in metrics. These problems become even more critical when downstream use cases depend on reliable data products, such as customer segmentation, churn analysis, recommendation systems, and marketing optimization.

A modern customer data platform must therefore combine multiple capabilities: it must ingest multi-format sources and event streams, store data efficiently and durably, support distributed processing, guarantee traceability and reproducibility, and provide a governance layer through consistent modeling and validation. From an engineering perspective, this leads naturally to a \textbf{lakehouse} architecture: a design that combines the scalability and flexibility of data lakes with the structure and reliability of data warehouses. In a lakehouse, raw data and curated data products coexist in a layered structure, while table formats and metadata services enforce consistency and enable reliable SQL access.

At the same time, business stakeholders expect reliable indicators that can be updated frequently and audited easily. This requires a data engineering foundation that is scalable, reproducible, and observable. In practice, such a foundation must address four core needs:
\begin{itemize}
    \item reliable ingestion of multi-source data, including event streams and batch files
    \item distributed storage with a clear organization strategy (raw, curated, analytics-ready)
    \item distributed processing and transformation with governance and data quality controls
    \item orchestration and monitoring to ensure operational reliability and traceability
\end{itemize}

This final year project presents \textbf{CustomerDNA AI -- Distributed Customer Data Lakehouse Platform}. The active Client~1 scope combines Bank Marketing, Online Shoppers Purchasing Intention, and Online Retail~II to represent campaign response, web-session intent, and transactional behavior. Kafka, HDFS, Spark, Iceberg, Hive Metastore, Trino, dbt-Spark, Great Expectations, Airflow, Prometheus, and Grafana form one governed and observable data engineering chain.

Beyond implementing the technical components, this project emphasizes operational engineering practices:
\begin{itemize}
    \item modularization of the platform into clear layers (ingestion, storage, processing, transformation, query, quality, orchestration, monitoring)
    \item a local control plane for orchestration and supervision, connected to a distributed data plane on three virtual machines
    \item idempotent reprocessing capabilities to support experimentation and debugging
    \item stored Great Expectations validation evidence and monitoring of service health and pipeline execution
\end{itemize}

The remainder of this report introduces the project context and Customer~360 objective, establishes the theoretical background, formalizes the requirements, justifies the selected technology stack, and then presents system design, implementation, distributed deployment, case-study validation, and performance evaluation. Future segmentation, churn, value prediction, personas, and recommendations are presented as downstream uses of the curated platform outputs.